% !TeX program = pdflatex
\documentclass[10pt,a4paper]{article}

\usepackage[utf8]{inputenc}
\usepackage[T1]{fontenc}
\usepackage{lmodern}
\usepackage{geometry}
\usepackage[hidelinks]{hyperref}
\usepackage{xcolor}

\geometry{
  top=1.7cm,
  bottom=1.8cm,
  left=1.8cm,
  right=1.8cm
}

\pagestyle{empty}
\frenchspacing
\setlength{\parindent}{0pt}
\setlength{\parskip}{0.45em}

\definecolor{mydarkgreen}{RGB}{0,0,0}

\begin{document}

\begin{minipage}[t]{0.48\textwidth}
\textbf{\large Javier Andres TARAZONA JIMENEZ}\\
Étudiant ingénieur IA \\
Machine Learning, Vision 3D, LLMs et prototypage \\
Massy, Île-de-France, France\\
\href{mailto:javitar06@gmail.com}{javitar06@gmail.com}\\
+33 07 46 36 97 75\\
\href{https://www.linkedin.com/in/javier-andres-tarazona-jimenez-84b489222/}{LinkedIn}
\end{minipage}
\hfill
\begin{minipage}[t]{0.42\textwidth}
\raggedleft
Safran Electronics \& Defense \\
Direction des Systèmes d’Information – Digital Factory \\
Massy, Île-de-France, France \\
\end{minipage}

\vspace{1.1cm}

\raggedleft
Palaiseau, le 24 juin 2026

\vspace{0.7cm}

\raggedright
\textbf{Candidature à l’alternance – Apprenti Ingénieur Data \& IA Cloud AWS}

\vspace{0.5cm}

{\small

Madame, Monsieur,

Actuellement étudiant colombien en double diplôme entre l’Universidad Nacional de Colombia et ENSTA Paris, au sein de l’Institut Polytechnique de Paris, je souhaite vous présenter ma candidature pour l’alternance d’Apprenti Ingénieur Data \& IA Cloud AWS au sein de Safran Electronics \& Defense, à partir de septembre 2026.

J’ai choisi l’informatique par intérêt pour la compréhension des systèmes complexes, par goût pour la construction de solutions concrètes à partir d’idées abstraites, et par volonté d’acquérir une autonomie technique solide. L’intelligence artificielle s’est ensuite imposée comme une évolution naturelle de ce parcours : passer de la programmation de systèmes à la conception de systèmes capables de percevoir, d’apprendre et d’aider à la décision. Ce domaine m’attire particulièrement parce qu’il combine rigueur mathématique, ambiguïté des problèmes ouverts et impact concret. La vision par ordinateur a d’abord connecté mon profil à une IA appliquée au monde réel ; aujourd’hui, l’IA générative, les LLMs et les architectures data modernes représentent pour moi une continuité stratégique de cette même ambition.

Safran m’intéresse précisément parce que l’IA n’y est pas seulement un outil d’optimisation numérique, mais une capacité opérationnelle appliquée à des systèmes industriels critiques. Votre positionnement dans l’aéronautique, le spatial et la défense, ainsi que la mission de Safran Electronics \& Defense autour des fonctions observer, décider et guider, correspondent fortement à ma vision de l’ingénierie : développer des technologies fiables, utiles et intégrées à des environnements où les contraintes de sécurité, de performance et de responsabilité sont réelles. Travailler sur l’IA et le cloud dans un tel contexte aurait pour moi une valeur particulière, car cela relie directement mes intérêts pour l’intelligence artificielle, les systèmes complexes, la mobilité, la défense et l’ingénierie appliquée.

L’offre proposée par la Digital Factory représente également une opportunité très cohérente avec mon parcours. Elle permet de réunir plusieurs dimensions que j’ai déjà explorées séparément : le développement logiciel, les pipelines de données, le NLP, les LLMs, les APIs, les bases de données et le déploiement applicatif. Participer à la conception d’une application complète permettant de déposer des documents, d’en extraire et structurer le contenu, de générer des embeddings, d’indexer les données dans une base vectorielle et d’interroger ces documents via un LLM correspond exactement au type de projet que je souhaite approfondir : une IA intégrée, industrialisable et utile aux métiers.

Je peux apporter à votre équipe un profil hybride entre intelligence artificielle, data engineering et développement logiciel. À ENSTA Paris – U2IS, je travaille sur des pipelines de données synthétiques, l’entraînement deep learning avec PyTorch et Vision Transformers, ainsi que la structuration de datasets pour l’apprentissage IA. Chez APEDYS91, j’ai développé un pipeline NLP de correction de texte en français avec spaCy, LanguageTool et des LLMs locaux, incluant des contrôles d’entités, de nombres et de cohérence. Chez Engin A.I, j’ai renforcé mon expérience en Python, architecture logicielle modulaire et optimisation de pipelines de traitement de données et d’images.

Mes projets personnels renforcent également cette adéquation. Avec ORIUN, j’ai conçu une plateforme full-stack avec Django REST Framework, PostgreSQL, Docker, Google Cloud APIs, authentification JWT, rapports et statistiques. Avec Sharp Sight, j’ai développé un backend FastAPI intégrant collecte automatisée, nettoyage, structuration de données et exposition d’endpoints REST. Ces expériences me donnent des bases concrètes pour contribuer à une application combinant interface web, API, traitement documentaire, données non structurées et architecture cloud.

Curieux, autonome et habitué à apprendre rapidement de nouvelles technologies, je serais particulièrement motivé à rejoindre la Digital Factory de Safran afin de contribuer à vos projets d’IA générative et de data platform, tout en consolidant mes compétences dans un environnement industriel exigeant, international et techniquement ambitieux.

Je reste à votre disposition pour un entretien afin de vous présenter plus en détail ma motivation et l’adéquation de mon profil avec cette alternance.

Veuillez agréer, Madame, Monsieur, l’expression de mes salutations distinguées.

\textbf{Cordialement},

Javier Andres TARAZONA JIMENEZ


}

\end{document}